\documentclass[UTF8,11pt]{ctexart}
\usepackage[a4paper,margin=24mm]{geometry}
\usepackage{amsmath,amssymb,amsthm,mathtools,mathrsfs}
\usepackage[all]{xy}
\usepackage{xurl,hyperref,enumitem}
\usepackage{tikz}
\usetikzlibrary{arrows.meta,cd,decorations.markings,calc,patterns}
\hypersetup{hidelinks,breaklinks=true}
\setlength{\emergencystretch}{3em}
\allowdisplaybreaks
\newtheorem*{theorem}{Theorem}
\newtheorem*{thm}{Theorem}
\newtheorem*{lemma}{Lemma}
\newtheorem*{proposition}{Proposition}
\newtheorem*{prop}{Proposition}
\newtheorem*{corollary}{Corollary}
\newtheorem*{cor}{Corollary}
\newtheorem*{claim}{Claim}
\theoremstyle{definition}
\newtheorem*{definition}{Definition}
\newtheorem*{defn}{Definition}
\newtheorem*{remark}{Remark}
\newtheorem*{example}{Example}
\newcommand{\R}{\mathbb R}
\newcommand{\C}{\mathbb C}
\newcommand{\Z}{\mathbb Z}
\newcommand{\Q}{\mathbb Q}
\newcommand{\N}{\mathbb N}
\newcommand{\F}{\mathbb F}
\newcommand{\D}{\mathbb D}
\newcommand{\bB}{\mathbb B}
\newcommand{\bC}{\mathbb C}
\newcommand{\bR}{\mathbb R}
\newcommand{\bZ}{\mathbb Z}
\newcommand{\bN}{\mathbb N}
\newcommand{\bQ}{\mathbb Q}
\newcommand{\bD}{\mathbb D}
\newcommand{\bF}{\mathbb F}
\newcommand{\CP}{\mathbb{CP}}
\newcommand{\RP}{\mathbb{RP}}
\newcommand{\sO}{\mathscr O}
\newcommand{\sI}{\mathscr I}
\newcommand{\sF}{\mathscr F}
\newcommand{\sG}{\mathscr G}
\newcommand{\sH}{\mathscr H}
\newcommand{\sR}{\mathscr R}
\newcommand{\sM}{\mathscr M}
\newcommand{\sV}{\mathscr V}
\newcommand{\sU}{\mathscr U}
\DeclareMathOperator{\Hom}{Hom}
\DeclareMathOperator{\Ext}{Ext}
\DeclareMathOperator{\Tor}{Tor}
\DeclareMathOperator{\Spec}{Spec}
\DeclareMathOperator{\Proj}{Proj}
\DeclareMathOperator{\supp}{supp}
\DeclareMathOperator{\im}{im}
\DeclareMathOperator{\id}{id}
\DeclareMathOperator{\Tr}{Tr}
\DeclareMathOperator{\tr}{tr}
\DeclareMathOperator{\rank}{rank}
\DeclareMathOperator{\ord}{ord}
\DeclareMathOperator{\dist}{dist}
\DeclareMathOperator{\End}{End}
\DeclareMathOperator{\Aut}{Aut}
\DeclareMathOperator{\Gal}{Gal}
\DeclareMathOperator{\vol}{vol}
\DeclareMathOperator{\Cl}{Cl}
\DeclareMathOperator{\coker}{coker}
\DeclareMathOperator{\codim}{codim}
\DeclareMathOperator{\divergence}{div}
\DeclareMathOperator{\Ric}{Ric}
\DeclareMathOperator{\grad}{grad}
\DeclareMathOperator{\Hess}{Hess}
\newcommand{\abs}[1]{\left\lvert#1\right\rvert}
\newcommand{\sabs}[1]{\lvert#1\rvert}
\newcommand{\babs}[1]{\bigl\lvert#1\bigr\rvert}
\newcommand{\Babs}[1]{\Bigl\lvert#1\Bigr\rvert}
\newcommand{\bbabs}[1]{\biggl\lvert#1\biggr\rvert}
\newcommand{\BBabs}[1]{\Biggl\lvert#1\Biggr\rvert}
\newcommand{\norm}[1]{\left\lVert#1\right\rVert}
\newcommand{\snorm}[1]{\lVert#1\rVert}
\newcommand{\bnorm}[1]{\bigl\lVert#1\bigr\rVert}
\newcommand{\Bnorm}[1]{\Bigl\lVert#1\Bigr\rVert}
\newcommand{\bbnorm}[1]{\biggl\lVert#1\biggr\rVert}
\newcommand{\BBnorm}[1]{\Biggl\lVert#1\Biggr\rVert}
\newcommand{\widebar}[1]{\overline{#1}}
\newcommand{\nicefrac}[2]{\frac{#1}{#2}}
\newcommand{\linnprod}[2]{\langle#1,#2\rangle}
\newcommand{\myindex}[1]{#1}
\newcommand{\glsadd}[1]{}
\newcommand{\avoidbreak}{}
\newcommand{\sourceref}[1]{\textnormal{[原文：\nolinkurl{#1}]}}
\newcommand{\thmref}[1]{\sourceref{#1}}
\newcommand{\lemmaref}[1]{\sourceref{#1}}
\newcommand{\propref}[1]{\sourceref{#1}}
\newcommand{\corref}[1]{\sourceref{#1}}
\newcommand{\defnref}[1]{\sourceref{#1}}
\newcommand{\exampleref}[1]{\sourceref{#1}}
\newcommand{\exerciseref}[1]{\sourceref{#1}}
\newcommand{\figureref}[1]{\sourceref{#1}}
\newcommand{\sectionref}[1]{\sourceref{#1}}
\newcommand{\appendixref}[1]{\sourceref{#1}}
\newcommand{\stacksref}[2]{\href{https://stacks.math.columbia.edu/tag/#1}{#2 [#1]}}

\begin{document}
\section*{046\quad Dedekind zeta函数的解析类数公式}
\subsection*{来源与计数目标}
Andrew V. Sutherland，MIT 18.785 \emph{Number Theory I}，Fall 2021，Lecture 19。11/15/2021；\S19.1.1 全部核心计数与体积推导（pp.~4--7），Theorem 19.8、Lemmas 19.9、19.11、Theorem 19.12 及其证明（pp.~7--9）；格点估计为 Corollary 19.6（p.~3）。采用 MIT OpenCourseWare 实际访问版本；获取日期：2026-09-28。
\par\url{https://ocw.mit.edu/courses/18-785-number-theory-i-fall-2021/mit18_785f21_lec19.pdf}
\par\textbf{本文件只计一个问题：}证明任意数域 Dedekind zeta 函数在 $\operatorname{Re}z>1-1/n$ 的亚纯延拓及 $z=1$ 留数公式，常数的体积与理想类来源均需展开。
\subsection*{作者命题与人类证明}
\paragraph{Definitions and the lattice-counting input.}
The Dedekind zeta function of a number field $K$ is defined by
\[
\zeta_K(s):=\sum_{\mathfrak a}N(\mathfrak a)^{-s}
=\prod_{\mathfrak p}(1-N(\mathfrak p)^{-s})^{-1},
\]
where $\mathfrak a$ ranges over nonzero $\mathcal O_K$-ideals and $\mathfrak p$ ranges over nonzero prime ideals. The product converges absolutely on $\operatorname{Re}(s)>1$ since $\#\{\mathfrak p\mid p\}\leq n:=[K:\Q]$ and $N(\mathfrak p):=[\mathcal O_K:\mathfrak p]\geq p$ imply that
\[
\sum_{\mathfrak p}|\log(1-N(\mathfrak p)^{-s})|\leq n\sum_p|\log(1-p^{-s})|,
\]
and the sum on the RHS converges on $\operatorname{Re}(s)>1$ since we know this holds for $\zeta_\Q(s)=\zeta(s)$.
\begin{definition}[19.4]
A set $B$ in a metric space $X$ is $d$-Lipschitz parametrizable if it is the union of the images of a finite number of Lipschitz continuous functions $f_i:[0,1]^d\to X$.
\end{definition}
\begin{corollary}[19.6; standard lattice-counting input]
Let $\Lambda$ be a lattice in an $\R$-vector space $V\simeq\R^n$ and let $S\subseteq V$ be a measurable set whose boundary is $(n-1)$-Lipschitz parametrizable. Then
\[
\#(tS\cap\Lambda)=\frac{\mu(S)}{\operatorname{covol}(\Lambda)}t^n+O(t^{n-1}).
\]
\end{corollary}
\paragraph{19.1.1 Counting algebraic integers of bounded norm}
Recall from \S15.2 that the unit group $K_\R^\times$ of $K_\R:=K\otimes_\Q\R$ is the locally compact group
\[
K_\R^\times\simeq\prod_{v\mid\infty}K_v^\times\simeq\prod_{\text{real }v\mid\infty}\R^\times\times\prod_{\text{complex }v\mid\infty}\C^\times.
\]
We have a natural embedding $K^\times\hookrightarrow K_\R^\times$, $x\mapsto(x_v)$, where $v$ ranges over the $r+s$ archimedean places of $K$; this allows us to view $K^\times$ as a subgroup of $K_\R^\times$ that contains the nonzero elements of $\mathcal O_K$. In Lecture 15 we defined the continuous homomorphism
\[
\operatorname{Log}:K_\R^\times\to\R^{r+s},\qquad (x_v)\mapsto(\log\|x_v\|_v),
\]
and proved that we have an exact sequence of abelian groups
\[
1\longrightarrow\mu_K\longrightarrow\mathcal O_K^\times\xrightarrow{\operatorname{Log}}\Lambda_K\longrightarrow0,
\]
in which $\Lambda_K$ is a lattice in the trace-zero hyperplane $\R^{r+s}_0:=\{x\in\R^{r+s}:T(x)=0\}$ (where $T(x)$ is the sum of the coordinates of $x$). The regulator $R_K$ is the covolume of $\Lambda_K$ in $\R^{r+s}_0$ (see Definition 15.16), where we endow $\R^{r+s}_0$ with the Euclidean measure induced by any coordinate projection $\R^{r+s}\to\R^{r+s-1}$. By Dirichlet's unit theorem (Theorem 15.12), we can write $\mathcal O_K^\times=U\times\mu_K$, where $U\subseteq\mathcal O_K^\times$ is free of rank $r+s-1$ (the subgroup $U$ is not uniquely determined, but let us fix a choice).

We want to estimate the quantity $\#\{\mathfrak a:N(\mathfrak a)\leq t\}$, where $\mathfrak a$ ranges over nonzero ideals of $\mathcal O_K$, as $t\to\infty$. As a first step, let us restrict our attention to principal ideals $(\alpha)\subseteq\mathcal O_K$. We then want to estimate the cardinality of $\{(\alpha):N(\alpha)\leq t\}$. For nonzero $\alpha,\alpha'\in K$ we have $(\alpha)=(\alpha')$ if and only if $\alpha/\alpha'\in\mathcal O_K^\times$, so this is equivalent to
\[
\{\alpha\in K^\times\cap\mathcal O_K:N(\alpha)\leq t\}/\mathcal O_K^\times,
\]
where for any set $S\subseteq K_\R^\times$, the notation $S/\mathcal O_K^\times$ denotes the set of equivalence classes of $S$ under the equivalence relation $\alpha\sim\alpha'\Leftrightarrow\alpha=u\alpha'$ for some $u\in\mathcal O_K^\times$. If we now define
\[
K_{\R,\leq t}^\times:=\{x\in K_\R^\times:N(x)\leq t\}\subseteq K_\R^\times\subseteq K_\R,
\]
then we want to estimate the cardinality of the finite set $(K_{\R,\leq t}^\times\cap\mathcal O_K)/\mathcal O_K^\times$, where the intersection takes place in $K_\R$ and produces a subset of $K_\R^\times$ that we partition into equivalence classes modulo $\mathcal O_K^\times$. To simplify matters, let us replace $\mathcal O_K^\times$ with the free group $U\subseteq\mathcal O_K^\times$; we then have a $w_K$-to-1 map
\[
(K_{\R,\leq t}^\times\cap\mathcal O_K)/U\longrightarrow(K_{\R,\leq t}^\times\cap\mathcal O_K)/\mathcal O_K^\times.
\]
It suffices to estimate the cardinality of $(K_{\R,\leq t}^\times\cap\mathcal O_K)/U$ and divide the result by $w_K$.

Recall that for $x=(x_v)\in K_\R^\times$, the norm map $N:K_\R^\times\to\R_{>0}^\times$ is defined by
\[
N(x):=\prod_{v\mid\infty}\|x_v\|_v=\prod_{v\text{ real}}|x_v|_\R\prod_{v\text{ complex}}|x_v|_\C^2,
\]
and satisfies $T(\operatorname{Log}x)=\log N(x)$ for all $x\in K_\R^\times$. We now define a surjective homomorphism
\[
\nu:K_\R^\times\twoheadrightarrow K_{\R,1}^\times,\qquad x\mapsto xN(x)^{-1/n}.
\]
The image of $K_{\R,1}^\times$ under the Log map is precisely the trace zero hyperplane $\R^{r+s}_0$ in $\R^{r+s}$ in which $\operatorname{Log}(U)=\operatorname{Log}(\mathcal O_K^\times)=\Lambda_K$ is a lattice. Let us fix a fundamental domain $F$ for the lattice $\Lambda_K$ in $\R^{r+s}_0$ so that
\[
S:=\nu^{-1}(\operatorname{Log}^{-1}(F))
\]
is a set of unique coset representatives for the quotient $K_\R^\times/U$. If we now define $S_{\leq t}:=\{x\in S:N(x)\leq t\}\subseteq K_\R$, we want to estimate the cardinality of the finite set $S_{\leq t}\cap\mathcal O_K$.

The set $\mathcal O_K$ is a lattice in the $\R$-vector space $K_\R$ of dimension $n$. We have $tS_{\leq1}=S_{\leq t^n}$, so we can estimate the cardinality of $S_{\leq t}=t^{1/n}S_{\leq1}$ via Corollary 19.6 with $S=S_{\leq1}$ and $\Lambda=\mathcal O_K$ by replacing $t$ with $t^{1/n}$, provided that the boundary of $S_{\leq1}$ is $(n-1)$-Lipschitz parametrizable, which we now argue.

The kernel of the Log map is $\{\pm1\}^r\times\mathrm U(1)^s$, where $\mathrm U(1):=\{z\in\C:z\bar z=1\}$ is the unit circle in $\C$. We thus have a continuous isomorphism of locally compact groups
\begin{align*}
K_\R^\times=(\R^\times)^r\times(\C^\times)^s&\xrightarrow{\sim}\R^{r+s}\times\{\pm1\}^r\times[0,2\pi)^s,\tag{1}\\
x=(x_1,\ldots,x_r,z_1,\ldots,z_s)&\longmapsto(\operatorname{Log}x)\times(\operatorname{sgn}x_1,\ldots,\operatorname{sgn}x_r)\times(\arg z_1,\ldots,\arg z_s),
\end{align*}
where the map to $\R^{r+s}$ is the Log map, the map to $\{\pm1\}^r$ is the vector of signs of the $r$ real components, and the map to $[0,2\pi)^s$ is the vector of angles $\arg z$ such that $z/|z|=e^{i\arg z}$ of the $s$ complex components.

The set $S_{\leq1}$ consists of $2^r$ connected components, one for each element of $\{\pm1\}^r$. We can parametrize each of these component using $n$ real parameters as follows:
\begin{itemize}
\item $r+s-1$ parameters in $[0,1)$ that encode a point in $F$ as an $\R$-linear combination of $\operatorname{Log}(\epsilon_1),\ldots,\operatorname{Log}(\epsilon_{r+s-1})$, where $\epsilon_1,\ldots,\epsilon_{r+s-1}$ are a basis for $U$,
\item $s$ parameters in $[0,1)$ that encode an element of $\mathrm U(1)^s$;
\item a parameter in $(0,1]$ that encodes the $n$th-root of the norm.
\end{itemize}
These parametrizations define a continuously differentiable bijection from the set $C=[0,1)^{n-1}\times(0,1]\subseteq[0,1]^n$ to each of the $2^r$ disjoint components of $S_{\leq1}$; it can be written out explicitly in terms of exponentials and the identity function. The boundary $\partial C$ is the boundary of the unit $n$-cube which is clearly $(n-1)$-Lipschitz parametrizable; thus each component of $S_{\leq1}$, and therefore $S_{\leq1}$ itself, has a boundary that is $(n-1)$-Lipschitz parametrizable.

We now apply Corollary 19.6 to the lattice $\mathcal O_K$ and the set $S_{\leq1}$ in the $n$-dimensional $\R$-vector space $K_\R$ with $t$ replaced by $t^{1/n}$, since $S_{\leq t}=t^{1/n}S_{\leq1}$. This yields
\[
\#(S_{\leq t}\cap\mathcal O_K)=\frac{\mu(S_{\leq1})}{\operatorname{covol}(\mathcal O_K)}(t^{1/n})^n+O((t^{1/n})^{n-1})
=\frac{\mu(S_{\leq1})}{|D_K|^{1/2}}t+O(t^{1-1/n}).\tag{2}
\]
Our next task is compute $\mu(S_{\leq1})$; as noted in Remark 19.7, we must use the normalized Haar measure $\mu$ on $K_\R$ defined in \S14.2 when doing so. We will use the isomorphism in (1) to make a change of coordinates, we just need to understand how this affects the Haar measure $\mu$ on $K_\R=\prod_{v\mid\infty}K_v\simeq\R^r\times\C^s$. In terms of the standard Lebesgue measures $dx$ and $dA$ on $\R$ and $\C$, we have $\mu=(dx)^r(2dA)^s$, where the $2dA$ reflects the fact that the normalized absolute value $\|\ \|_v$ for each complex place $v$ is the square of the Euclidean absolute value on $\C$. For each factor of $K_\R^\times=\prod_{v\mid\infty}K_v\simeq(\R^\times)^r\times(\C^\times)^s\subseteq\R^r\times\C^s$ we define the maps
\begin{align*}
\R^\times&\longrightarrow\R\times\{\pm1\},&\C^\times&\longrightarrow\R\times[0,2\pi),\\
x&\longmapsto(\log|x|,\operatorname{sgn}x),&z&\longmapsto(2\log|z|,\arg z),\\
\pm e^\ell&\longleftarrow(\ell,\pm1),&e^{\ell/2+i\theta}&\longleftarrow(\ell,\theta),\\
dx&\longmapsto e^\ell d\ell\mu_{\{\pm1\}},&2dA&\longmapsto2e^{\ell/2}d(e^{\ell/2})d\theta=e^\ell d\ell d\theta,
\end{align*}
where $d\ell$ is the Lebesgue measure on $\R$, $\mu_{\{\pm1\}}$ is the counting measure on $\{\pm1\}$, and $d\theta$ is the Lebesgue measure on $[0,2\pi)$. We thus have
\[
K_\R^\times\xrightarrow{\sim}\R^{r+s}\times\{\pm1\}^r\times[0,2\pi)^s,\qquad
\mu\mapsto e^{T(\cdot)}\mu_{\R^{r+s}}\mu_{\{\pm1\}}^r\mu_{[0,2\pi)}^s,
\]
where the trace function $T(\cdot)$ sums the coordinates of a vector in $\R^{r+s}$. We now make one further change of coordinates:
\begin{align*}
\R^{r+s}&\longrightarrow\R^{r+s-1}\times\R,\\
x=(x_1,\ldots,x_{r+s})&\longmapsto(x_1,\ldots,x_{r+s-1},y:=T(x)),\\
e^{T(x)}\mu_{\R^{r+s}}&\longmapsto e^y\mu_{\R^{r+s-1}}dy.
\end{align*}
If we let $\pi_0:\R^{r+s}\to\R^{r+s-1}$ denote the coordinate projection, then the Lebesgue measure of $\pi_0(F)$ in $\R^{r+s-1}$ is, by definition, the regulator $R_K$ (see Definition 15.16). The Log map gives us a bijection
\begin{align*}
S_{\leq1}&\xrightarrow{\sim}F+(-\infty,0]\left(\tfrac1n,\ldots,\tfrac1n,\tfrac2n,\ldots,\tfrac2n\right),\\
x=N(x)^{1/n}\nu(x)&\longmapsto\operatorname{Log}\nu(x)+\log N(x)\left(\tfrac1n,\ldots,\tfrac1n,\tfrac2n,\ldots,\tfrac2n\right).
\end{align*}
The coordinate $y\in(-\infty,0]$ is given by $y=T(\operatorname{Log}x)=\log N(x)$, so we can view $S_{\leq1}$ as an infinite union of cosets of $\operatorname{Log}^{-1}(F)$ parameterized by $e^y=N(x)\in(0,1]$. Under our change of coordinates we thus have
\[
\begin{aligned}
K_\R^\times&\xrightarrow{\sim}\R^{r+s-1}\times\R\times\{\pm1\}^r\times[0,2\pi)^s,\\
S_{\leq1}&\longrightarrow\pi_0(F)\times(-\infty,0]\times\{\pm1\}^r\times[0,2\pi)^s.
\end{aligned}
\]
Since $R_K=\mu_{\R^{r+s-1}}(\pi(F))$, we have
\[
\mu(S_{\leq1})=\int_{-\infty}^0 e^y R_K2^r(2\pi)^s\,dy=2^r(2\pi)^sR_K.
\]
Plugging this into (2) yields
\[
\#(S_{\leq t}\cap\mathcal O_K)=\frac{2^r(2\pi)^sR_K}{|D_K|^{1/2}}t+O(t^{1-1/n}).\tag{3}
\]
\paragraph{19.2 Proof of the analytic class number formula}
\begin{theorem}[19.8]
Let $K$ be a number field of degree $n$. As $t\to\infty$, the number of nonzero $\mathcal O_K$-ideals $\mathfrak a$ of absolute norm $N(\mathfrak a)\leq t$ is
\[
\frac{2^r(2\pi)^sh_KR_K}{w_K|D_K|^{1/2}}t+O(t^{1-1/n}),
\]
where $r$ and $s$ are the number of real and complex places of $K$, respectively, $h_K:=\#\operatorname{cl}\mathcal O_K$ is the class number, $R_K$ is the regulator, $w_K:=\#\mu_K$ is the number of roots of unity, and $D_K:=\operatorname{disc}\mathcal O_K$ is the absolute discriminant.
\end{theorem}
\begin{proof}
In order to count the nonzero $\mathcal O_K$-ideals $\mathfrak a$ of absolute norm $N(\mathfrak a)\leq t$ we group them by ideal class. For the trivial class, we just need to count nonzero principal ideals $(\alpha)$, equivalently, the number of nonzero $\alpha\in\mathcal O_K$ with $N(\alpha)\leq t$, modulo the unit group $\mathcal O_K^\times$. Dividing (3) by $w_K$ to account for the $w_K$-to-1 map
\[
S_{\leq t}\cap\mathcal O_K\longrightarrow(K_{\R,\leq t}^\times\cap\mathcal O_K)/\mathcal O_K^\times,
\]
we obtain
\[
\#\{(\alpha)\subseteq\mathcal O_K:N(\alpha)\leq t\}=\frac{2^r(2\pi)^sR_K}{w_K|D_K|^{1/2}}t+O(t^{1-1/n}).\tag{4}
\]
To complete the proof we now show that we get the same answer for every ideal class; the nonzero ideals $\mathfrak a$ of norm $N(\mathfrak a)\leq t$ are asymptotically equidistributed among ideal classes.

Fix an ideal class $[\mathfrak a]$, with $\mathfrak a\subseteq\mathcal O_K$ nonzero (every ideal class contains an integral ideal, by Theorem 14.20). Multiplication by $\mathfrak a$ gives a bijection
\begin{align*}
\{\text{ideals }\mathfrak b\in[\mathfrak a^{-1}]:N(\mathfrak b)\leq t\}
&\xrightarrow{\times\mathfrak a}\{\text{nonzero principal ideals }(\alpha)\subseteq\mathfrak a:N(\alpha)\leq tN(\mathfrak a)\}\\
&\longrightarrow\{\text{nonzero }\alpha\in\mathfrak a:N(\alpha)\leq tN(\mathfrak a)\}/\mathcal O_K^\times.
\end{align*}
Let $S_{[\mathfrak a],\leq t}$ denote the set on the RHS. The estimate in (4) derived from Corollary 19.6 applies to any lattice in $K_\R$, not just $\mathcal O_K$. Replacing $\mathcal O_K$ with $\mathfrak a$ in (4) we obtain
\begin{align*}
\#S_{[\mathfrak a],\leq t}
&=\frac{2^r(2\pi)^sR_K}{w_K\operatorname{covol}(\mathfrak a)}tN(\mathfrak a)+O(t^{1-1/n})\\
&=\frac{2^r(2\pi)^sR_K}{w_K\operatorname{covol}(\mathcal O_K)N(\mathfrak a)}tN(\mathfrak a)+O(t^{1-1/n})\\
&=\frac{2^r(2\pi)^sR_K}{w_K|D_K|^{1/2}}t+O(t^{1-1/n}),
\end{align*}
since $\operatorname{covol}(\mathfrak a)=N(\mathfrak a)\operatorname{covol}(\mathcal O_K)$, by Corollary 14.17. Note that the RHS does not depend on the ideal class $[\mathfrak a]$. Summing over ideal classes yields
\[
\#\{\text{nonzero ideals }\mathfrak b\subseteq\mathcal O_K:N(\mathfrak b)\leq t\}
=\sum_{[\mathfrak a]\in\operatorname{cl}(\mathcal O_K)}\#S_{[\mathfrak a],\leq t}
=\frac{2^r(2\pi)^sh_KR_K}{w_K|D_K|^{1/2}}t+O(t^{1-1/n}),
\]
as claimed.
\end{proof}
\begin{lemma}[19.9]
Let $a_1,a_2,\ldots$ be a sequence of complex numbers and let $\sigma$ be a real number. Suppose that $a_1+\cdots+a_t=O(t^\sigma)$ as $t\to\infty$. Then the Dirichlet series $\sum a_nn^{-s}$ defines a holomorphic function on $\operatorname{Re}s>\sigma$.
\end{lemma}
\begin{proof}
Let $A(x):=\sum_{0<n\leq x}a_n$. Writing the Dirichlet sum as a Stieltjes integral (apply Corollary 18.28 with $f(n)=n^{-s}$ and $g(n)=a_n$), for $\operatorname{Re}(s)>\sigma$ we have
\begin{align*}
\sum_{n=1}^\infty a_nn^{-s}
&=\int_{1^-}^\infty x^{-s}\,dA(x)
=\left.\frac{A(x)}{x^s}\right|_{1^-}^\infty-\int_{1^-}^\infty A(x)\,dx^{-s}\\
&=(0-0)-\int_{1^-}^\infty A(x)(-sx^{-s-1})\,dx
=s\int_{1^-}^\infty\frac{A(x)}{x^{s+1}}\,dx.
\end{align*}
Note that we used $|A(x)|=O(x^\sigma)$ and $\operatorname{Re}(s)>\sigma$ to conclude $\lim_{x\to\infty}A(x)/x^s=0$. The integral on the RHS converges locally uniformly on $\operatorname{Re}(s)>\sigma$ and the lemma follows.
\end{proof}
\begin{lemma}[19.11]
Let $a_1,a_2,\ldots$ be a sequence of complex numbers that satisfies $a_1+\cdots+a_t=\rho t+O(t^\sigma)$ as $t\to\infty$ for some $\sigma\in[0,1)$ and $\rho\in\C^\times$. The Dirichlet series $\sum a_nn^{-s}$ converges on $\operatorname{Re}(s)>1$ and has a meromorphic continuation to $\operatorname{Re}(s)>\sigma$ that is holomorphic except for a simple pole at $s=1$ with residue $\rho$.
\end{lemma}
\begin{proof}
Define $b_n:=a_n-\rho$. Then $b_1+\cdots+b_t=O(t^\sigma)$ and
\[
\sum a_nn^{-s}=\rho\sum n^{-s}+\sum b_nn^{-s}=\rho\zeta(s)+\sum b_nn^{-s}.
\]
We have already proved that the Riemann zeta function $\zeta(s)$ is holomorphic on $\operatorname{Re}(s)>1$ and has a meromorphic continuation to $\operatorname{Re}(s)>0$ that is holomorphic except for a simple pole at 1 with residue 1. By the previous lemma, $\sum b_nn^{-s}$ is holomorphic on $\operatorname{Re}(s)>\sigma$, and since $\sigma<1$, it is holomorphic at $s=1$. So the entire RHS has a meromorphic continuation to $\operatorname{Re}(s)>\sigma$ that is holomorphic except for the simple pole at 1 coming from $\zeta(s)$, and the residue at $s=1$ is $\rho\cdot1+0=\rho$.
\end{proof}
\begin{theorem}[19.12 (Analytic Class Number Formula)]
Let $K$ be a number field of degree $n$. The Dedekind zeta function $\zeta_K(z)$ extends to a meromorphic function on $\operatorname{Re}(z)>1-\frac1n$ that is holomorphic except for a simple pole at $z=1$ with residue
\[
\lim_{z\to1^+}(z-1)\zeta_K(z)=\rho_K:=\frac{2^r(2\pi)^sh_KR_K}{w_K|D_K|^{1/2}},
\]
where $r$ and $s$ are the number of real and complex places of $K$, respectively, $h_K:=\#\operatorname{cl}\mathcal O_K$ is the class number, $R_K$ is the regulator, $w_K:=\#\mu_K$ is the number of roots of unity, and $D_K:=\operatorname{disc}\mathcal O_K$ is the absolute discriminant.
\end{theorem}
\begin{proof}
We have $\zeta_K(z)=\sum_{\mathfrak a}N(\mathfrak a)^{-z}=\sum_{t\geq1}a_tt^{-z}$, where $\mathfrak a$ ranges over nonzero ideals of $\mathcal O_K$, and $a_t:=\#\{\mathfrak a:N(\mathfrak a)=t\}$ with $t\in\Z_{\geq1}$. If we now define $\rho_K:=2^r(2\pi)^sh_KR_K/(w_K|D_K|^{1/2})$, then by Theorem 19.8 we have
\[
a_1+\cdots+a_t=\#\{\mathfrak a:N(\mathfrak a)\leq t\}=\rho_Kt+O(t^{1-1/n})\quad(\text{as }t\to\infty).
\]
Applying Lemma 19.11 with $\sigma=1-1/n$, we see that $\zeta_K(z)=\sum a_tt^{-z}$ extends to a meromorphic function on $\operatorname{Re}(z)>1-1/n$ that is holomorphic except for a simple pole at $z=1$ with residue $\rho_K$.
\end{proof}
\subsection*{整理说明（非作者正文）}
开头从第1页摘取定义，Corollary 19.6 明列为标准格点估计先修；未声称包含它的证明。其余核心段从 19.1.1 的单位商、区域构造、边界参数化、测度换元，到理想类等分布和解析留数均完整收录。辅助估计、19.8 和19.12 不拆开计数。原文中的 $[0,2\pi)$ 群坐标、一次 $\pi$/$\pi_0$ 写法、Log 所诱导映射被称为 bijection 等表述照录，未新增其证明或静默改写。

标准先修为数域整数环及理想的格嵌入与余体积公式、Dirichlet 单位定理、类群有限性（本库002、010有对应人类证明）、Corollary 19.6 的 Lipschitz 边界格点估计、Stieltjes 分部积分和 Riemann zeta 在1处的简单极点。完整的单位商区域和留数常数没有当作先修。第1、5--9页已实际查看图像，第2--4页仅取得解析文本，未声称全部图像核对。难度依据是把乘法单位格与加法理想格同时组织为带精确常数的计数，再转为解析极点；不是重复002的有限性或010的单位结构命题。
\subsection*{可修改的简短标注（非作者正文）}
$c$：数域类数与 zeta 留数

$a$：单位群的对数格；理想格余体积；调节子；Lipschitz 边界格点估计；基本区域；实复位测度换元；理想类等分布；Stieltjes 分部积分

\texttt{cross\_theory\_status = yes}。单位群商的几何基本区域把调节子和无穷位因子变成体积；理想格计数给出类数控制的线性主项，再由 Dirichlet 级数分部积分将其变成 zeta 留数。位置：\S19.1.1、19.8、19.11--19.12。

全部标注均为非穷尽、可修改初稿；未标注不表示未使用。
\subsection*{许可与修改记录}
材料来自 MIT OpenCourseWare，作者与课程见来源。按该站所示 CC BY-NC-SA 4.0 许可复用，整理部分沿用同一许可。\url{https://creativecommons.org/licenses/by-nc-sa/4.0/}。2026-09-28：助手转录题证、调整数学排版并加入来源、范围和初稿标注；不暗示作者或 MIT 认可本次整理。所留为本次题证转录摘录，不是原 PDF 下载件。
\end{document}
